% =============================================================================
% abstract_methods_draft.tex
% Paper-agent draft — gap depth study, 1 M_Jup at 20 au
% Generated: 2026-05-29  |  Model: Claude Sonnet 4.6 (GitHub Copilot)
% Task ID: paper_agent_gap_depth_20260529
%
% Reference bibcodes verified via NASA ADS (mcp_ads_ads_search, 2026-05-29).
% References marked [VERIFY IN ADS] could not be confirmed at draft time.
%
% BibTeX key → ADS bibcode mapping (all verified unless noted):
%   Kanagawa2016      → 2016PASJ...68...43K
%   Kanagawa2017      → 2017PASJ...69...97K  [VERIFY IN ADS]
%   Andrews2018       → 2018ApJ...869L..41A
%   Huang2018         → 2018ApJ...869L..42H
%   Zhang2018         → 2018ApJ...869L..47Z
%   Dullemond2018     → 2018ApJ...869L..46D
%   Crida2006         → 2006Icar..181..587C
%   ShakuraSunyaev1973 → 1973A&A....24..337S
%   Stammler2022      → 2022ApJ...935...35S
%   BenitezLlambay2016 → 2016ApJS..223...11B
%   Masset2000        → 2000A&AS..141..165M
%   Birnstiel2010     → 2010A&A...513A..79B
%   Birnstiel2012     → 2012A&A...539A.148B
%   deValBorro2006    → 2006MNRAS.370..529D
% =============================================================================

% ---------------------------------------------------------------------------
% ABSTRACT
% ---------------------------------------------------------------------------

\begin{abstract}
Planet-carved gaps in protoplanetary discs are powerful tracers of embedded
giant planets and are now routinely resolved at millimetre wavelengths by
the Atacama Large Millimetre/submillimetre Array (ALMA), as demonstrated by
the Disk Substructures at High Angular Resolution Project
\citep[DSHARP;][]{Andrews2018, Huang2018}.
We present a combined analytical and numerical investigation of the gap
opened by a $1\,M_\mathrm{Jup}$ planet orbiting a solar-mass star at
$a_\mathrm{p} = 20\,\mathrm{au}$ in a disc characterised by a constant
aspect ratio $h/r = 0.05$ and a Shakura-Sunyaev viscosity parameter
$\alpha = 10^{-3}$ \citep{ShakuraSunyaev1973}.
The dimensionless Kanagawa parameter $K \equiv q^2/(h^5\alpha) = 2916$
places the planet deep in the non-linear gap-opening regime, yielding an
analytical prediction of $\Sigma_\mathrm{gap}/\Sigma_0 = 8.5\times10^{-3}$
(118-fold gas depletion; \citealt{Kanagawa2016}).
Complementary 1D dust-evolution models computed with \textsc{DustPy}
\citep{Stammler2022} reveal a fragmentation-limited grain population at
$r = 10\,\mathrm{au}$ ($a_\mathrm{max} \simeq 2\times10^{-4}\,\mathrm{cm}$)
and a drift-limited regime at $r = 100\,\mathrm{au}$
($a_\mathrm{max} \simeq 24\,\mathrm{cm}$) after $10^3\,\mathrm{yr}$ of
evolution with $v_\mathrm{frag} = 10\,\mathrm{m\,s}^{-1}$.
Full 2D hydrodynamical simulations with \textsc{FARGO3D}
\citep{BenitezLlambay2016} are underway to provide a direct numerical test
of the \citet{Kanagawa2016} prediction.
\end{abstract}

% ---------------------------------------------------------------------------
% METHODS SECTION
% ---------------------------------------------------------------------------

\section{Methods}
\label{sec:methods}

\subsection{Hydrodynamical simulations (\textsc{FARGO3D})}
\label{sec:fargo}

We perform two-dimensional isothermal hydrodynamical simulations of the
gas disc using \textsc{FARGO3D} \citep{BenitezLlambay2016}, a
GPU-oriented, multi-fluid disc code that employs the orbital-advection
algorithm of \citet{Masset2000} to remove the dominant azimuthal velocity
component before the Courant integration, allowing time-steps orders of
magnitude larger than in a standard Eulerian scheme.

The computational domain spans $r \in [6, 60]\,\mathrm{au}$, corresponding
to $[0.3, 3.0]$ in units of the planet semi-major axis
$a_\mathrm{p} = 20\,\mathrm{au}$, and is discretised on a polar grid of
$N_r \times N_\phi = 256 \times 512$ logarithmically spaced radial and
uniformly spaced azimuthal cells.
The initial surface-density profile follows a power law,
$\Sigma(r) = \Sigma_0\,(r/a_\mathrm{p})^{-0.5}$, with normalisation
$\Sigma_0 = 6.8\times10^{-4}\,M_\star\,a_\mathrm{p}^{-2}$.
The disc aspect ratio is held constant at $h/r = 0.05$ (zero flaring index),
and the \citet{ShakuraSunyaev1973} $\alpha$-viscosity prescription is
applied with $\alpha = 10^{-3}$.

A single planet of mass $M_\mathrm{p} = 1\,M_\mathrm{Jup}$ is placed on a
fixed circular orbit at $a_\mathrm{p}$; its gravitational potential is
softened over $0.6\,H$ to avoid numerical singularities.
Wave-killing zones following \citet{deValBorro2006} are applied inside
$r < 0.3\times1.15\,a_\mathrm{p}$ and outside
$r > 3.0/1.15\,a_\mathrm{p}$ to damp spurious reflected waves at the
radial boundaries.
The simulation is integrated for 1000 planetary orbits
($t \approx 89\,\mathrm{kyr}$), which is sufficient for the gap to reach a
quasi-steady state \citep{Kanagawa2016}.

\subsection{Dust evolution (\textsc{DustPy})}
\label{sec:dustpy}

We model the 1D radial evolution of the dust size distribution using
\textsc{DustPy} \citep{Stammler2022}, a Python implementation of the
coagulation--fragmentation--drift framework of
\citet{Birnstiel2010}.
Grain growth, fragmentation, and radial drift are solved self-consistently
on 80 logarithmically spaced radial cells spanning $r = 1$--$100\,\mathrm{au}$.
The stellar parameters are $M_\star = 1\,M_\odot$ and
$T_\star = 4000\,\mathrm{K}$; the total disc mass is
$M_\mathrm{disc} = 0.05\,M_\odot$.
We adopt the same turbulence parameter $\alpha = 10^{-3}$ used in the
hydrodynamical simulations and a fragmentation velocity
$v_\mathrm{frag} = 10\,\mathrm{m\,s}^{-1}$, consistent with
laboratory measurements for silicate aggregates \citep{Birnstiel2012}.
In the inner disc ($r \lesssim 30\,\mathrm{au}$), where the sound speed is
elevated, turbulence-driven relative velocities set the grain-size ceiling
(fragmentation barrier); in the cold outer disc grains instead grow until
radial drift removes them before the next collision \citep{Birnstiel2012}.
The dust simulation is run for $t_\mathrm{end} = 10^3\,\mathrm{yr}$, and
the maximum grain size $a_\mathrm{max}(r)$ is recorded at the final snapshot.

\subsection{Analytical benchmark: Kanagawa gap-depth formula}
\label{sec:analytical}

As an independent benchmark we evaluate the empirical gap-depth formula
derived by \citet{Kanagawa2016} from a suite of 2D isothermal simulations:
%
\begin{equation}
  \frac{\Sigma_\mathrm{gap}}{\Sigma_0} = \frac{1}{1 + 0.04\,K},
  \qquad
  K \equiv \frac{q^2}{h^5\,\alpha},
  \label{eq:kanagawa_depth}
\end{equation}
%
where $q = M_\mathrm{p}/M_\star$ is the planet-to-star mass ratio.
For our fiducial parameters ($q = 9.55\times10^{-4}$, $h = 0.05$,
$\alpha = 10^{-3}$) we obtain $K = 2916$ and
$\Sigma_\mathrm{gap}/\Sigma_0 = 8.5\times10^{-3}$, a 118-fold depletion.
The gap-opening condition is evaluated via two complementary criteria: the
thermal criterion $M_\mathrm{p} > 3\,M_\star h^3$
($M_\mathrm{p,th} = 0.39\,M_\mathrm{Jup}$; \citealt{Crida2006}) and the
viscous Kanagawa criterion $K > K_\mathrm{open} = 25$
($M_\mathrm{p,visc} = 0.09\,M_\mathrm{Jup}$; \citealt{Kanagawa2016}).
Both thresholds are exceeded by factors of $2.5\times$ and $117\times$,
respectively, confirming that the planet is strongly non-linear.
The \citet{Kanagawa2016} formula was calibrated over the range
$K \lesssim \mathcal{O}(10^3)$; at $K = 2916$ we assign a conservative
systematic uncertainty of $\pm30\%$ on $\Sigma_\mathrm{gap}/\Sigma_0$
\citep{Kanagawa2017}.
Planet masses inferred from gap widths following \citet{Zhang2018} and the
broader DSHARP gap statistics of \citet{Huang2018} will serve as the
observational benchmark for the completed simulation.

% ---------------------------------------------------------------------------
% END OF DRAFT
% ---------------------------------------------------------------------------
